\documentclass[10pt,a4paper]{amsart}
\usepackage[margin=19mm]{geometry}
\usepackage{amsmath,amssymb,mathtools}
\usepackage[scr=rsfs,cal=euler]{mathalfa}
\usepackage{xfrac}
\usepackage[unicode,hidelinks,bookmarks=false]{hyperref}
\newcommand{\ie}{i.e.\ }
\newcommand{\cf}{cf.\ }
\newcommand{\resp}{respectively\ }
\newcommand{\Subsection}{\subsection}
\providecommand{\nobreakdash}{\nobreak\hbox{-}}
\setlength{\emergencystretch}{2em}
\multlinegap0pt
\usepackage{amssymb}
\usepackage[shortlabels]{enumitem}
\newtheorem{theorem}{Theorem}
\title{A 920-block explicit construction guaranteeing a triple intersection with every $6$-subset of $[60]$}
\author{Paulo Henrique Cunha Gomes}
\date{Source: arXiv:2601.06179v1 (CC BY 4.0)}
\begin{document}
\maketitle

\noindent\emph{Source and licence.} A 920-block explicit construction guaranteeing a triple intersection with every $6$-subset of $[60]$. Version arXiv:2601.06179v1 (CC BY 4.0), distributed by arXiv under the Creative Commons Attribution 4.0 International licence, http://creativecommons.org/licenses/by/4.0/, as recorded in the arXiv licence metadata for this exact version and shown on the abstract page https://arxiv.org/abs/2601.06179v1. Full licence notice: \texttt{licenses/arxiv-2601.06179-CC-BY-4.0.txt}.\par\smallskip

\noindent\textit{Editorial scope.} Counted unit: the article's theorem thm:main (source0.tex lines 79--83). The article's complete original proof of it is delivered with it (source0.tex lines 85--101, one \texttt{proof} environment of 17 lines). No mathematics is written, repaired or re-derived: the statement and the proof are the article's own lines, in the article's own notation, and no retained line is substituted or re-flowed.  No further source-local context is needed: every cross-reference of every retained line resolves inside the two delivered spans.  Nothing was omitted from any retained span, so the package carries no omission record.  No retained line contains a citation, so the wrapper carries no bibliography and no bibliography entry.  No retained line contains a figure environment, an image-inclusion command or drawing code, so no image is delivered and the package declares no asset dependency.  Editorial formatting only, in addition: an AMS \texttt{amsart} preamble that restates the article's own declarations and macros used by the retained lines, namely the environments \texttt{theorem} on separate counters, together with the standard packages those lines need.  The retained environments print in this document as Theorem 1 in order of appearance, whereas the article prints the same items with the numbering of its own full document.  The complete native source of the article as distributed in the arXiv e-print for this exact version is bundled unchanged as \texttt{sources/192394/source0.tex}.
\begin{theorem}\label{thm:main}
Let $\mathcal{B}$ be the family of $920$ blocks constructed above.
For every $S\subset[60]$ with $|S|=6$, there exists $B\in\mathcal{B}$ such that
$|S\cap B|\ge 3$.
\end{theorem}

\begin{proof}
Let $S\subset[60]$ with $|S|=6$, and write $S_A=S\cap A$ and $S_C=S\cap C$.
Since $|S_A|+|S_C|=6$, either $|S_A|\ge 3$ or $|S_C|\ge 3$.

Assume $|S_A|\ge 3$ (the case $|S_C|\ge 3$ is identical with $Q$-pairs).
Choose distinct elements $a,b,c\in S_A$.
Each of $a,b,c$ belongs to a unique pair among $\{P_1,\dots,P_{15}\}$; denote these pairs by
$P_{i_a},P_{i_b},P_{i_c}$.

If $i_a,i_b,i_c$ are distinct, then
$B=P_{i_a}\cup P_{i_b}\cup P_{i_c}\in\mathcal{B}_A$ contains $\{a,b,c\}$.

If two of the indices coincide, say $P_{i_a}=P_{i_b}$, choose any $r\notin\{i_a,i_c\}$ and set
$B=P_{i_a}\cup P_{i_c}\cup P_r\in\mathcal{B}_A$.
Again $B$ contains $\{a,b,c\}$.
Thus in all cases $|S\cap B|\ge 3$.
\end{proof}

\end{document}
